\section{Land Cover Data Cleaning and Feature Preparation}

\subsection{Classification Code Harmonization}

The raw $\texttt{lcccode}$ column contains complex codes (e.g., \texttt{7001 // 8001}) which represent mixed, uncertain, or transitional land cover types. The separator (\texttt{//}) indicates a mixture between the two classes. To handle this complexity for machine learning, the following steps were taken:

\textbf{Legend Integration}
The official legend files (\texttt{globcover\_LCCS\_legend\_africa.xls} and \texttt{globcover\_legend.xls}) were loaded to establish a mapping between the numeric codes and descriptive labels.

\begin{itemize}
    \item The $\texttt{LCCCode}$ was mapped to its descriptive label ($\texttt{LCCOwnLabel}$) using a dictionary.
    \item This process identified that certain codes were designated as ``Error'' in the $\texttt{LC}$ column of the legend, yet represented valid land cover types in the $\texttt{LCCOwnLabel}$ column (e.g., ``Closed forest''). This confirmed that these codes should be retained and mapped correctly.
\end{itemize}

\textbf{High-Level Mapping}
Given the 22 unique $\texttt{lcccode}$ values in the dataset, a high-level, simplified mapping was created to aggregate these codes into fundamental categories (e.g., ``Forests,'' ``Bare lands,'' ``Croplands,'' ``Vegetation,'' etc.). 
For instance:
\begin{itemize}
    \item Multiple codes corresponding to water-related classes (e.g., 7001, 8001) were grouped as ``Water bodies''.
    \item Forest-related codes (e.g., 21497-121340) were grouped as ``Forests''.
\end{itemize}
This high-level categorization reduces the number of classes and improves model interpretability, while preserving the main land cover distinctions.

\begin{table}[H]
    \centering
    \caption{Example of Landcover Transformation from \texttt{lcccode} to \texttt{lcc\_highlevel}}
    \label{tab:landcover_highlevel}
    \begin{tabular}{cc}
        \toprule
        \textbf{lcccode} & \textbf{lcc\_highlevel} \\
        \midrule
        7001 // 8001      & Water bodies \\
        0011      & Bare Lands \\
        11498     & Croplands \\
        21497-121340      & Forests \\
        11490 // 11494      & Croplands \\
        \bottomrule
    \end{tabular}
\end{table}
